% Code Quest – Algorithms World – Level 1 Issue 1
% Această fișă LaTeX reprezintă structura completă a primului număr
% „CALCULATORUL NU ŞTIE SĂ NUMERE”. Include atât text narativ cât şi
% exerciții de diverse tipuri. Poți compila acest fișier cu pdflatex.

\documentclass[11pt,a4paper]{article}

\usepackage[utf8]{inputenc}
\usepackage[T1]{fontenc}
\usepackage[romanian]{babel}
\usepackage{geometry}
\geometry{margin=1in}
\usepackage{xcolor}
\usepackage{tcolorbox}
\usepackage{listings}
\usepackage{amsmath,amssymb}
\usepackage{graphicx}
\usepackage{float}

% Definim stilul pentru cod C++
\lstdefinestyle{cpp}{
  language=C++,
  basicstyle=\ttfamily\footnotesize,
  keywordstyle=\color{blue}\bfseries,
  commentstyle=\color{gray}\itshape,
  numbers=left,
  numberstyle=\tiny,
  stepnumber=1,
  numbersep=5pt,
  showspaces=false,
  showstringspaces=false,
  showtabs=false,
  frame=single,
  tabsize=2
}

% Cutii personalizate pentru diverse blocuri pedagogice
\newtcolorbox{missionbox}[1][]{title=#1, colback=blue!5, colframe=blue!40, boxrule=1pt, arc=2mm, fonttitle=\bfseries}
\newtcolorbox{tryit}[1][]{title=#1, colback=green!5, colframe=green!40, boxrule=1pt, arc=2mm, fonttitle=\bfseries}
\newtcolorbox{examplebox}[1][]{title=#1, colback=yellow!5, colframe=yellow!40, boxrule=1pt, arc=2mm, fonttitle=\bfseries}

\title{Code Quest -- Lumea Algoritmilor\\Nivelul 1 -- Edi\c{t}ia 1}
\author{}
\date{}

\begin{document}

\maketitle

\begin{center}
  {\Large \textbf{Calculatorul nu \c{s}tie s\u{a} numere}}\\[0.5em]
  {\small O lec\c{t}ie de la variabile p\^an\u{a} la p\u{a}trate}
\end{center}

\section*{Introducere}
Calculatorul nu \c{s}tie s\u{a} numere f\u{a}r\u{a} ajutorul programatorului. El \c{s}tie doar s\u{a} \c{t}in\u{a} minte numere \c{s}i s\u{a} fac\u{a} aritmetic\u{a} simpl\u{a}. Pentru a-l \^{i}nv\u{a}\c{t}a s\u{a} numere, vom descoperi variabilele, operatorii aritmetici, structurile de control \c{s}i modul \^{i}n care putem construi forme pe un ecran, exact ca \^{i}n Minecraft.

\begin{missionbox}[Misiunea ta]
\textbf{Scopul acestei edi\c{t}ii}: vei \^{i}nv\u{a}\c{t}a s\u{a} programezi un calculator s\u{a} numere \c{s}i s\u{a} deseneze p\u{a}trate folosind bucle. Vei porni de la no\c{t}iuni de baz\u{a} despre variabile \c{s}i vei ajunge la dublu \texttt{for} \c{s}i diagone. La final, vei fi capabil s\u{a} construie\c{s}ti un render de p\u{a}trat \^{i}ntr-un sistem de coordonate.
\end{missionbox}

\section*{Memoria \c{s}i variabilele}
Variabilele sunt zone de memorie \^{i}n care calculatoarele stocheaz\u{a} date. Cel mai des utilizat tip este \texttt{int}, care reprezint\u{a} un num\u{a}r \^{i}ntreg (pozitiv sau negativ, f\u{a}r\u{a} zecimale).

\begin{examplebox}[Exemplu de declarare de variabile]
\begin{lstlisting}[style=cpp]
int a = 1;
int b = 10;
int i = 0;
\end{lstlisting}
\end{examplebox}

\begin{tryit}[\^{I}ncearc\u{a} singur]
Declar\u{a} trei variabile \^{i}ntregi: \texttt{x}, \texttt{y} \c{s}i \texttt{scor} \c{s}i ini\c{t}ializeaz\u{a}-le cu valorile \texttt{5}, \texttt{12} \c{s}i \texttt{0}. Scrie codul t\u{a}u mai jos:
\end{tryit}

\section*{Operatori aritmetici}
Folosim operatori aritmetici pentru a modifica valorile variabilelor.

\begin{examplebox}[Incrementare \c{s}i decrementare]
\begin{lstlisting}[style=cpp]
i = i + 1;   // cre\c{s}te i cu 1
i += 1;      // cre\c{s}te i cu 1
i++;         // cre\c{s}te i cu 1

i = i - 1;   // scade i cu 1
i -= 1;      // scade i cu 1
i--;         // scade i cu 1
\end{lstlisting}
\end{examplebox}

\begin{tryit}[Exerci\c{t}iu rapid]
Dac\u{a} \texttt{i = 7}, care va fi valoarea lui \texttt{i} dup\u{a} fiecare dintre instruc\c{t}iunile de mai jos?
\begin{enumerate}
\item \texttt{i++;}
\item \texttt{i += 3;}
\item \texttt{i--;}
\end{enumerate}
\end{tryit}

\section*{Structuri de control}
Calculatorul poate "s\u{a}r\u{a}" la alte instruc\c{t}iuni, ceea ce permite controlul fluxului programului. Cea mai simpl\u{a} structur\u{a} de control este condi\c{t}ionala \texttt{if}.

\begin{examplebox}[Structura \texttt{if}]
\begin{lstlisting}[style=cpp]
if(i <= 10) {
    cout << i;
}
\end{lstlisting}
Instruc\c{t}iunea \texttt{cout << i;} se execut\u{a} doar dac\u{a} condi\c{t}ia \texttt{i <= 10} este adev\u{a}rat\u{a}.
\end{examplebox}

\begin{tryit}[Condi\c{t}ii]
Scrie o instruc\c{t}iune \texttt{if} care verific\u{a} dac\u{a} variabila \texttt{x} este pozitiv\u{a}. Dac\u{a} este, afiseaz\u{a} mesajul "num\u{a}r pozitiv".
\end{tryit}

\section*{Bucle \texttt{while}}
O bucl\u{a} permite repetarea instruc\c{t}iunilor c\u{a}t timp o condi\c{t}ie este adev\u{a}rat\u{a}.

\begin{examplebox}[Buclele \texttt{while}]
\begin{lstlisting}[style=cpp]
int i = 1;
while(i <= 10) {
    cout << i << " ";
    i++;
}
\end{lstlisting}
\end{examplebox}

\begin{tryit}[Num\u{a}rare cu \texttt{while}]
Modific\u{a} programul de mai sus astfel \^{i}nc\u{a}t s\u{a} afiseze numerele de la \texttt{1} la \texttt{20}.
\end{tryit}

\section*{Cum num\u{a}r\u{a}m?}
Pentru a programa calculatorul s\u{a} numere, trebuie s\u{a} definim trei valori:
\begin{enumerate}
\item De unde s\u{a} \^{i}nceap\u{a} – \texttt{int i = 1;}
\item C\u{a}t timp s\u{a} numere – \texttt{while(i <= 10)}
\item Cu ce pas s\u{a} numere – \texttt{i++;}
\end{enumerate}

\section*{Bucle \texttt{for}}
Structura \texttt{for} este o form\u{a} compact\u{a} de bucl\u{a}.

\begin{examplebox}[Structura \texttt{for}]
\begin{lstlisting}[style=cpp]
for(int i = 1; i <= 10; i++) {
    cout << i << " ";
}
\end{lstlisting}
Aceasta echivaleaz\u{a} cu:
\begin{lstlisting}[style=cpp]
int i = 1;
while(i <= 10) {
    cout << i << " ";
    i++;
}
\end{lstlisting}
\end{examplebox}

\begin{tryit}[Parcurgere cu \texttt{for}]
Scrie un program care afiseaz\u{a} numerele pare de la \texttt{2} la \texttt{20} folosind o bucl\u{a} \texttt{for}.
\end{tryit}

\section*{Bucl\u{a} \texttt{for} combinat\u{a} cu \texttt{if}}
Putem combina structuri pentru a realiza comportamente mai complexe. Un exemplu clasic este FizzBuzz:

\begin{examplebox}[FizzBuzz]
\begin{lstlisting}[style=cpp]
for(int i = 1; i <= n; i++) {
    if(i % 3 == 0 && i % 5 == 0) {
        cout << "FizzBuzz";
    } else if(i % 3 == 0) {
        cout << "Fizz";
    } else if(i % 5 == 0) {
        cout << "Buzz";
    } else {
        cout << i;
    }
    cout << " ";
}
\end{lstlisting}
\end{examplebox}

\begin{tryit}[Exerci\c{t}iu FizzBuzz]
Modific\u{a} programul \texttt{for} astfel \^{i}nc\u{a}t s\u{a} afiseze:
\begin{itemize}
\item "Fizz" pentru multiplii de 3;
\item "Buzz" pentru multiplii de 5;
\item "FizzBuzz" pentru multiplii de 3 \c{s}i 5;
\item altfel num\u{a}rul \^{i}n sine.
\end{itemize}
\end{tryit}

\section*{Bucle \^{i}mbricate: dublu \texttt{for}}
Pentru a desena forme 2D (cum ar fi un p\u{a}trat) avem nevoie de dou\u{a} bucle: una pentru r\u{a}nduri \c{s}i una pentru coloane.

\begin{examplebox}[P\u{a}trat simplu]
\begin{lstlisting}[style=cpp]
for(int linie = 1; linie <= 5; linie++) {
    for(int coloana = 1; coloana <= 5; coloana++) {
        cout << "#";
    }
    cout << endl;
}
\end{lstlisting}
Output-ul este:
\begin{verbatim}
#####
#####
#####
#####
#####
\end{verbatim}
\end{examplebox}

\begin{tryit}[Deseneaz\u{a} un p\u{a}trat]
Scrie un program care deseneaz\u{a} un p\u{a}trat de dimensiune \texttt{n x n} folosind simbolul \texttt{*}. Alege tu dimensiunea \texttt{n}.
\end{tryit}

\section*{Diagonalele unui p\u{a}trat}
Pentru un p\u{a}trat de m\u{a}rime \texttt{n x n}, elementele aflate pe diagonala principal\u{a} \^{i}ndeplinesc condi\c{t}ia \texttt{linie == coloana}. Pe diagonala secundar\u{a}, condi\c{t}ia este \texttt{linie + coloana == n + 1}.

\begin{tryit}[Diagona Principal\u{a}]
Scrie un program care afiseaz\u{a} un p\u{a}trat \texttt{n x n} \^{i}n care doar elementele de pe diagonala principal\u{a} sunt \texttt{*}, iar restul \texttt{.}.
\end{tryit}

\begin{tryit}[Diagona secundar\u{a}]
Scrie un program care afiseaz\u{a} un p\u{a}trat \texttt{n x n} \^{i}n care doar elementele de pe diagonala secundar\u{a} sunt \texttt{*}, iar restul \texttt{.}.
\end{tryit}

\section*{Forma final\u{a}: p\u{a}trat cu margini \c{s}i diagonale}
Combin\u{a} condi\c{t}iile pentru a desena un p\u{a}trat care are:
\begin{itemize}
\item \# pe margini;
\item * pe diagonala principal\u{a};
\item + pe diagonala secundar\u{a};
\item . pe restul.
\end{itemize}

\begin{tryit}[Problema final\u{a}]
Se cite\c{s}te \texttt{n}. Deseneaz\u{a} p\u{a}tratul \texttt{n x n} urm\u{a}nd regulile de mai sus. Aminte\c{s}te-\c{t}i priorit\u{a}\c{t}ile: mai \^{i}nt\u{a}i marginea, apoi diagonalele.
\end{tryit}

\section*{Concluzie}
Prin aceast\u{a} lec\c{t}ie am \^{i}nv\u{a}\c{t}at cum putem \^{i}nv\u{a}\c{t}a calculatorul s\u{a} numere \c{s}i s\u{a} deseneze forme prin folosirea variabilelor, operatorilor aritmetici, structurilor de control \texttt{if}, buclelor \texttt{while} \c{s}i \texttt{for}. Aceste concepte stau la baza program\u{a}rii de jocuri, graficii \c{s}i multor algoritmi.

\end{document}